%% ----------------------------------------------------------------------------
%%  §3.7  Successioni monotòne
%% ----------------------------------------------------------------------------
\section{Successioni monotòne}
\label{sec:monotone}

Abbiamo già dato nel \cref{sec:invertibili} la definizione di funzioni (strettamente)
crescenti o decrescenti. Allo stesso modo, per le successioni diciamo che:

\begin{definizione}[Successioni monotòne]\index{successione!monotòna}
Una successione $a_{n}$ si dice
\begin{align}
  &\text{\emph{strettamente crescente}} \text{ se } & a_{n}&<a_{n+1} & &\perogni n\in\N; \label{eq:succ-str-cresc}\\
  &\text{\emph{crescente}} \text{ se } & a_{n}&\le a_{n+1} & &\perogni n\in\N; \label{eq:succ-cresc}\\
  &\text{\emph{strettamente decrescente}} \text{ se } & a_{n}&>a_{n+1} & &\perogni n\in\N; \label{eq:succ-str-decr}\\
  &\text{\emph{decrescente}} \text{ se } & a_{n}&\ge a_{n+1} & &\perogni n\in\N. \label{eq:succ-decr}
\end{align}
Una successione $a_{n}$ è \emph{monotòna} se verifica una delle quattro condizioni sopra
scritte; è \emph{strettamente monotòna} se verifica la \eqref{eq:succ-str-cresc} oppure la
\eqref{eq:succ-str-decr}.
\end{definizione}

Basta confrontare ogni termine con il successivo: per induzione, se $a_{n}\le a_{n+1}$
per ogni $n$, allora $a_{n}\le a_{m}$ per ogni $n\le m$, e la successione è crescente
anche nel senso della definizione data per le funzioni. Se $a_{n}=a$ per ogni $n\in\N$,
con $a$ numero reale fissato, si dice che $a_{n}$ è una successione \emph{costante}. Le
successioni costanti sono allo stesso tempo crescenti e decrescenti.

\begin{esempio}
Le successioni \eqref{eq:succ-1n} e \eqref{eq:succ-n1n} sono strettamente monotòne: la
\eqref{eq:succ-1n} $a_{n}=1/n$ è strettamente decrescente, mentre la \eqref{eq:succ-n1n}
$a_{n}=(n-1)/n$ è strettamente crescente; infatti, se $a_{n}=1/n$, allora
\begin{equation}\label{eq:es-monotona-1}
  a_{n}>a_{n+1} \quad\sse\quad \frac1n>\frac1{n+1} \quad\sse\quad n+1>n,
\end{equation}
e l'ultima disuguaglianza è verificata per ogni $n\in\N$. Invece, se $a_{n}=(n-1)/n$,
\begin{equation}\label{eq:es-monotona-2}
  a_{n}<a_{n+1} \quad\sse\quad \frac{n-1}{n}<\frac{n}{n+1} \quad\sse\quad n^{2}-1<n^{2},
\end{equation}
e anche questa volta l'ultima disuguaglianza è verificata per ogni $n\in\N$.
\end{esempio}

Il risultato seguente è di fondamentale importanza: nello schema della
\cref{fig:schema-teoremi} è il primo dei grandi teoremi di esistenza che seguono
l'esistenza dell'estremo superiore, e da essa discende direttamente.

\begin{teorema}[sulle successioni monotòne]\label{teo:monotone}\index{teorema!sulle successioni monotòne}
Ogni successione monotòna ammette limite. In particolare, ogni successione monotòna e
limitata è convergente, cioè ammette limite finito. Più precisamente:
\begin{equation}\label{eq:lim-monotone}
\begin{aligned}
  a_{n}\ \text{crescente} &\implica \lim_{n\to+\infty}a_{n}=\sup_{n\in\N}a_{n};\\
  a_{n}\ \text{decrescente} &\implica \lim_{n\to+\infty}a_{n}=\inf_{n\in\N}a_{n}.
\end{aligned}
\end{equation}
\end{teorema}

\begin{schema}
\Ipotesi{$a_{n}$ monotòna (ad esempio crescente).}
\Tesi{$a_{n}$ è regolare, e $\lim a_{n}=\sup_{n}a_{n}$ (che è finito se e solo se $a_{n}$
è limitata superiormente, e vale $+\infty$ altrimenti).}
\end{schema}

\begin{strategia}
Il candidato limite di una successione crescente è il «tetto» dei suoi valori, cioè
$\ell=\sup_{n}a_{n}$. Per dimostrare che è il limite servono due cose, e ciascuna viene da
uno dei due ingredienti:
\begin{itemize}
  \item i termini non superano mai $\ell$, perché $\ell$ è un maggiorante; quindi
  $a_{n}<\ell+\epsilon$ per ogni $n$, gratis;
  \item i termini si avvicinano a $\ell$ \emph{da sotto}: per la caratterizzazione
  dell'estremo superiore (\cref{prop:caratt-sup}) c'è \emph{almeno un} termine $a_{\nu}$
  maggiore di $\ell-\epsilon$; ed essendo la successione crescente, \emph{tutti} i
  termini successivi lo sono a maggior ragione.
\end{itemize}
È la monotonia che trasforma «almeno un termine» in «definitivamente tutti»
(\cref{fig:monotona-sup}): senza di essa la successione potrebbe avvicinarsi a $\ell$ e poi allontanarsene. Se la successione
non è limitata superiormente, lo stesso ragionamento con un livello $M$ al posto di
$\ell-\epsilon$ dà la divergenza a $+\infty$. Le successioni decrescenti si trattano in
modo simmetrico, con l'estremo inferiore.
\end{strategia}

\begin{dimostrazione}
Consideriamo il caso di una successione $a_{n}$ crescente e limitata superiormente.
Posto $\ell=\sup_{n}a_{n}$, che esiste ed è finito per il \cref{teo:esistenza-sup},
fissato $\epsilon>0$, per le proprietà dell'estremo superiore
(\cref{prop:caratt-sup}) esiste $\nu\in\N$ tale che
\begin{equation}\label{eq:monotone-passo1}
  \ell-\epsilon<a_{\nu}.
\end{equation}
Per $n>\nu$ risulta $a_{\nu}\le a_{n}$, e d'altra parte $a_{n}\le\ell$ perché $\ell$ è un
maggiorante; dunque
\begin{equation}\label{eq:monotone-passo2}
  \ell-\epsilon<a_{\nu}\le a_{n}\le\ell<\ell+\epsilon \qquad\perogni n>\nu,
\end{equation}
da cui $\lim_{n\to+\infty}a_{n}=\ell$.

Consideriamo ora il caso di una successione $a_{n}$ crescente e non limitata
(superiormente). Fissato $M>0$, poiché $M$ non è un maggiorante dei valori, esiste
$\nu\in\N$ tale che $a_{\nu}>M$. Dato che $a_{n}$ è crescente, per ogni $n>\nu$ risulta
\begin{equation}\label{eq:monotone-passo3}
  a_{n}\ge a_{\nu}>M,
\end{equation}
da cui $\lim_{n\to+\infty}a_{n}=+\infty=\sup_{n}a_{n}$.

In modo analogo si trattano i casi relativi a successioni decrescenti; oppure ci si
riconduce ai precedenti osservando che, se $a_{n}$ è decrescente, la successione $-a_{n}$
è crescente, e che $\inf_{n}a_{n}=-\sup_{n}(-a_{n})$.
\end{dimostrazione}

\begin{figure}[htb]
\centering
\begin{tikzpicture}
\begin{axis}[width=0.92\textwidth, height=5cm, xmin=0, xmax=21, ymin=0, ymax=2.35,
  xtick={5,10,15,20}, extra x ticks={7}, extra x tick labels={$\nu$},
  ytick={2}, yticklabels={$\ell=\sup a_{n}$}, xlabel={$n$}]
  \fill[verde!15] (axis cs:0,1.7) rectangle (axis cs:21,2.3);
  \addplot[tratteggio, domain=0:21] {2};
  \addplot[only marks, mark=*, mark size=1.4pt, blu, samples at={1,...,20}] {2-1.6*0.78^x};
  \addplot[tratteggio] coordinates {(7,0) (7,1.7)};
  \node[font=\scriptsize, verde!50!black, anchor=east] at (axis cs:20.8,1.78) {$\ell-\epsilon$};
\end{axis}
\end{tikzpicture}
\caption{Una successione crescente e limitata: appena un termine $a_{\nu}$ supera
$\ell-\epsilon$, la monotonia costringe tutti i termini successivi a restare nella
striscia $(\ell-\epsilon,\ell]$.}
\label{fig:monotona-sup}
\end{figure}

\begin{osservazione}
La successione \eqref{eq:succ-alt-n} $a_{n}=(-1)^{n}/n$ non è monotòna; infatti i termini di
posto pari sono positivi, mentre quelli di posto dispari sono negativi: i termini della
successione oscillano intorno allo zero. La successione $(-1)^{n}/n$ è quindi un
\emph{esempio di successione convergente pur non essendo monotòna} (osserviamo che ciò non
contraddice il teorema: infatti nel teorema non si afferma che ogni successione
convergente è monotòna!). La successione non regolare \eqref{eq:succ-alt} $a_{n}=(-1)^{n}$
non è monotòna; ciò è in accordo con il teorema sulle successioni monotòne, perché se la
successione $(-1)^{n}$ fosse monotòna dovrebbe avere limite.
\end{osservazione}

Ricordando che una successione si dice regolare se essa ammette limite (finito o
infinito), il teorema precedente afferma che \emph{ogni successione monotòna è regolare}.
Due osservazioni tecniche ne estendono l'uso.

\begin{osservazione}[Successioni definitivamente monotòne e successioni traslate]\label{oss:traslata}
\leavevmode
\begin{enumerate}
  \item Il teorema vale anche per le successioni \emph{definitivamente} monotòne, cioè
  monotòne a partire da un certo indice $n_{0}$: basta applicarlo alla successione
  $a_{n_{0}},a_{n_{0}+1},\dots$, perché i primi termini non influiscono sul limite. In tal
  caso, però, il limite non è più necessariamente l'estremo superiore di tutti i termini.
  \item Se $a_{n}\to a$ (finito o infinito), allora anche la successione «traslata»
  $a_{n+1}$ tende ad $a$: se $\abs{a_{n}-a}<\epsilon$ per $n>\nu$, a maggior ragione
  $\abs{a_{n+1}-a}<\epsilon$ per $n>\nu$. Lo stesso vale per $a_{n+k}$, con $k\in\N$ fisso.
\end{enumerate}
\end{osservazione}

\subsection{Un'applicazione: il metodo di Erone per le radici quadrate}

Il teorema sulle successioni monotòne è lo strumento naturale per studiare le
\emph{successioni definite per ricorrenza}\index{successione!definita per ricorrenza},
in cui ogni termine è calcolato a partire dal precedente. È così che funzionano quasi
tutti gli algoritmi numerici iterativi.

\begin{esempio}[Il metodo di Erone]\label{es:erone}\index{metodo di Erone}
Consideriamo la successione definita da
\begin{equation}\label{eq:erone}
  x_{1}=2, \qquad x_{n+1}=\frac12\Bigl(x_{n}+\frac2{x_{n}}\Bigr) \quad (n\in\N).
\end{equation}
L'idea è semplice: se $x_{n}$ è un'approssimazione di $\sqrt2$ per eccesso, allora
$2/x_{n}$ lo è per difetto (il loro prodotto vale $2$), e la loro media è
un'approssimazione migliore. Dimostriamo che $x_{n}\to\sqrt2$.

\emph{Tutti i termini sono positivi}: per induzione, se $x_{n}>0$ anche $x_{n+1}>0$.

\emph{Per $n\ge2$ si ha $x_{n}\ge\sqrt2$}: infatti, per ogni $n$,
\[
  x_{n+1}-\sqrt2=\frac{x_{n}^{2}-2\sqrt2\,x_{n}+2}{2x_{n}}=\frac{(x_{n}-\sqrt2)^{2}}{2x_{n}}\ge0
\]
(ed è $x_{1}=2\ge\sqrt2$ anche per $n=1$).

\emph{La successione è decrescente}: per ogni $n$,
\[
  x_{n+1}-x_{n}=\frac{2-x_{n}^{2}}{2x_{n}}\le0,
\]
perché $x_{n}\ge\sqrt2$, cioè $x_{n}^{2}\ge2$.

Per il \cref{teo:monotone}, essendo decrescente e limitata inferiormente da $\sqrt2$, la
successione converge a un limite $L$, e per il \cref{cor:confronto-limiti} $L\ge\sqrt2>0$.
Passiamo al limite nella \eqref{eq:erone}: il primo membro $x_{n+1}$ tende a $L$
(\cref{oss:traslata}), e il secondo, per il \cref{teo:operazioni-limiti}, tende a
$\frac12\bigl(L+\frac2L\bigr)$. Per l'unicità del limite
$L=\frac12\bigl(L+\frac2L\bigr)$, cioè $L^{2}=2$, e quindi $L=\sqrt2$.
\end{esempio}

\begin{attenzione}[Prima l'esistenza, poi il valore]
Il passaggio al limite nella relazione di ricorrenza è lecito \emph{solo dopo} aver
dimostrato che il limite esiste. Senza questa precauzione si possono ottenere assurdità:
per la successione $x_{1}=1$, $x_{n+1}=2x_{n}$ (cioè $x_{n}=2^{n-1}$), «passando al
limite» si otterrebbe $L=2L$, cioè $L=0$; ma la successione diverge a $+\infty$.
\end{attenzione}

\begin{python}[La convergenza del metodo di Erone]
from decimal import Decimal, getcontext
getcontext().prec = 60              # lavoriamo con 60 cifre decimali

radice_vera = Decimal(2).sqrt()
x = Decimal(2)
for n in range(1, 8):
    errore = x - radice_vera
    print(f"x_{n} = {str(x)[:34]}   errore = {float(errore):.1e}")
    x = (x + 2 / x) / 2
\end{python}
\begin{risultato}
x_1 = 2   errore = 5.9e-01
x_2 = 1.5   errore = 8.6e-02
x_3 = 1.41666666666666666666666666666666   errore = 2.5e-03
x_4 = 1.41421568627450980392156862745098   errore = 2.1e-06
x_5 = 1.41421356237468991062629557889013   errore = 1.6e-12
x_6 = 1.41421356237309504880168962350253   errore = 9.0e-25
x_7 = 1.41421356237309504880168872420969   errore = 2.9e-49
\end{risultato}

Il programma mostra un fenomeno notevole: a ogni passo il numero di cifre esatte
\emph{raddoppia} (l'errore passa da $10^{-6}$ a $10^{-12}$, poi a $10^{-24}$). È la
cosiddetta \emph{convergenza quadratica}, che si legge direttamente nella formula
$x_{n+1}-\sqrt2=\frac{(x_{n}-\sqrt2)^{2}}{2x_{n}}$: l'errore al passo successivo è
proporzionale al \emph{quadrato} dell'errore attuale (si veda anche la \cref{fig:erone}). Il
metodo di Erone è un caso
particolare del metodo di Newton per la ricerca degli zeri, che studieremo con le derivate
e che è alla base dei solutori non lineari usati nel calcolo strutturale e nella
fluidodinamica computazionale.

\begin{figure}[htb]
\centering
\includegraphics[width=0.62\textwidth]{erone}
\caption{Il diagramma «a ragnatela» del metodo di Erone. Per rendere visibili più passi
si parte da $x_{1}=3$: dal punto $(x_{1},x_{1})$ della bisettrice $y=x$ si scende
verticalmente alla curva $y=\frac12(x+\frac2x)$, trovando $x_{2}$, poi ci si sposta
orizzontalmente fino alla bisettrice, e così via. Le iterate scendono verso il punto
fisso $\sqrt2$, dove la curva incontra la bisettrice (\texttt{python/fig\_erone.py}).}
\label{fig:erone}
\end{figure}

\begin{notastorica}[Le radici quadrate nell'antichità]
Il metodo porta il nome di Erone di Alessandria (I secolo d.C.), che lo descrive nei
\emph{Metrica} per calcolare $\sqrt{720}$; ma era noto molto prima. La tavoletta
babilonese YBC 7289, conservata a Yale e risalente a circa il 1800 a.C., mostra un
quadrato con le diagonali e, scritto in base $60$, il numero $1;24,51,10$, cioè
$1+\frac{24}{60}+\frac{51}{60^{2}}+\frac{10}{60^{3}}\approx1{,}41421296$: un'approssimazione
di $\sqrt2=1{,}41421356\ldots$ con un errore di circa $6\cdot10^{-7}$ (cinque decimali
esatti). Curiosamente, essa coincide con il troncamento a tre cifre sessagesimali di
$x_{4}=577/408=1;24,51,10,35\ldots$, il quarto termine del metodo di Erone a partire da
$x_{1}=2$. Secondo alcuni storici i Babilonesi conoscevano già un procedimento iterativo
di questo tipo, che sarebbe allora il più antico algoritmo della storia.
\end{notastorica}
